\section{Introduction}
\label{sec:intro}
% BUDGET: about 3/4 of column 1.  Four paragraphs.
%
% P1  What the field measures now: Dice, clDice \cite{shit2021cldice},
%     Betti errors, cbDice \cite{shi2024cbdice}.
% P2  Why a run-length metric is different: it has a UNIT.  "You can trace
%     1052 px before hitting an error" is a sentence a clinician can check.
%     ERL comes from connectomics \cite{januszewski2018ffn}; state plainly
%     that it is NOT ours.
% P3  Two topology-metric surveys \cite{berger2024pitfalls, decroocq2025bench}
%     do not list it -- so it is unused in this field, which is the opening.
% P4  Contributions, as a list.  Be careful here: \cite{berger2024pitfalls}
%     already published "conventions distort the ranking" on DRIVE, so the
%     contribution is that it happens on ERL, along ERL's own axes.

\TODO{introduction}
% Budget: ~285 words of prose (3/4 column minus the contributions list).
\FILL{2-4}

\noindent\textbf{Contributions.}
\begin{enumerate}
  \item A constructed case set on which Dice ties to four decimals, clDice
        ties inside the window a reader would call equivalent, and Betti-0
        \emph{ranks the worse prediction better}; ERL separates them.
  \item The first execution of the reference ERL implementation on retinal
        data, which exposes an undocumented connectivity assumption and a
        fragment length that depends on input order.
  \item A specification: three under-specified axes, two further conditions,
        and the rule that governs how much the choice costs.
\end{enumerate}
